% SOURCE: Chinese PDF viewer pp.170--173, including the short final page 173.
\subsection*{3) Single-particle continuous spectrum}
For a particle-number representation $A$ with a discrete spectrum, the annihilation and creation operators are $\op a_i,\op a_i^\dagger$. We have
\[\op N=\sum_i\op n_i,\qquad \op n_i=\op a_i^\dagger\op a_i,\qquad
\braket{a_i}{a_j}=\delta_{ij},\qquad \sum_i\ket{a_i}\bra{a_i}=\mathbf1.\]
The commutation relations are
\begin{align*}
\text{Bosons:}&&
[\op a_i,\op a_j]&=0,&
[\op a_i^\dagger,\op a_j^\dagger]&=0,&
[\op a_i,\op a_j^\dagger]&=\delta_{ij},\\
\text{Fermions:}&&
\{\op a_i,\op a_j\}&=0,&
\{\op a_i^\dagger,\op a_j^\dagger\}&=0,&
\{\op a_i,\op a_j^\dagger\}&=\delta_{ij}.
\end{align*}
Now extend this to the particle-number representation for a continuous spectrum, for example the coordinate and spin states $\ket{\vect r,\sigma}$ (continuous coordinates and discrete spin). The annihilation and creation operators are $\op\psi_\sigma(\vect r)$ and $\op\psi_\sigma^\dagger(\vect r)$. Their basis satisfies
\[\op{\vect r}\ket{\vect r,\sigma}=\vect r\ket{\vect r,\sigma},\qquad
\op\sigma\ket{\vect r,\sigma}=\sigma\ket{\vect r,\sigma},\]
\[\braket{\vect r',\sigma'}{\vect r'',\sigma''}
=\delta^{(3)}(\vect r'-\vect r'')\delta_{\sigma'\sigma''},\qquad
\int\dd^3r\sum_\sigma\ket{\vect r,\sigma}\bra{\vect r,\sigma}=\mathbf1.\]
The particle-number operator is
\[\op N=\int\dd^3r\sum_\sigma\op n_\sigma(\vect r),\qquad
\op n_\sigma(\vect r)=\op\psi_\sigma^\dagger(\vect r)\op\psi_\sigma(\vect r).\]
The commutation relations become
\begin{align*}
\text{Bosons:}\quad
[\op\psi_{\sigma'}(\vect r'),\op\psi_{\sigma''}(\vect r'')]&=0,\\
[\op\psi_{\sigma'}^\dagger(\vect r'),\op\psi_{\sigma''}^\dagger(\vect r'')]&=0,\\
[\op\psi_{\sigma'}(\vect r'),\op\psi_{\sigma''}^\dagger(\vect r'')]
&=\delta^{(3)}(\vect r'-\vect r'')\delta_{\sigma'\sigma''};
\end{align*}
\begin{align*}
\text{Fermions:}\quad
\{\op\psi_{\sigma'}(\vect r'),\op\psi_{\sigma''}(\vect r'')\}&=0,\\
\{\op\psi_{\sigma'}^\dagger(\vect r'),\op\psi_{\sigma''}^\dagger(\vect r'')\}&=0,\\
\{\op\psi_{\sigma'}(\vect r'),\op\psi_{\sigma''}^\dagger(\vect r'')\}
&=\delta^{(3)}(\vect r'-\vect r'')\delta_{\sigma'\sigma''}.
\end{align*}
The additive one-particle operator has the form
\begin{align*}
\op{\mathcal A}&=\sum_{j,k}\op b_j^\dagger\op b_k\mel{b_j}{\op A}{b_k}\\
&\longrightarrow\sum_{\sigma',\sigma''}\int\dd^3r'\,\dd^3r''\,
\op\psi_{\sigma'}^\dagger(\vect r')\op\psi_{\sigma''}(\vect r'')
\mel{\vect r',\sigma'}{\op A}{\vect r'',\sigma''}.
\end{align*}
An additive two-particle operator (a two-particle potential) has the form
\begin{align*}
\op{\mathcal V}
&=\frac12\sum_{k,l,m,n}\op b_k^\dagger\op b_l^\dagger\op b_m\op b_n
\mel{b_k,b_l}{\op V}{b_m,b_n}\\
&\longrightarrow\frac12\sum_{\sigma_1,\sigma_2,\sigma_3,\sigma_4}
\int\dd^3r_1\,\dd^3r_2\,\dd^3r_3\,\dd^3r_4\,
\op\psi_{\sigma_1}^\dagger(\vect r_1)
\op\psi_{\sigma_2}^\dagger(\vect r_2)
\op\psi_{\sigma_3}(\vect r_3)
\op\psi_{\sigma_4}(\vect r_4)\\
&\hspace{2.7em}\times
\mel{\vect r_1,\sigma_1;\vect r_2,\sigma_2}{\op V}
{\vect r_3,\sigma_3;\vect r_4,\sigma_4}.
\end{align*}

\subsection*{4) Quantum dynamics and second quantization}
In the Heisenberg picture, consider the equation of motion for the annihilation operator:
\[\ii\hbar\frac{\dd\op a_i}{\dd t}=[\op a_i,\op H].\]
Neglect interactions involving three or more particles, retaining only one-particle and two-particle contributions:
\[\op H=\op H_0+\op V.\]
Here $\op H_0$ contains the kinetic energy of one particle and its potential energy in an external field. In an arbitrary discrete particle-number representation $A$,
\begin{align*}
\op H
&=\sum_{k,l}\op a_k^\dagger\op a_l\mel{a_k}{\op H_0}{a_l}\\
&\quad+\frac12\sum_{k,l,m,n}
\op a_k^\dagger\op a_l^\dagger\op a_m\op a_n
\mel{a_k,a_l}{\op V}{a_m,a_n}.
\end{align*}
Since
\begin{align*}
[\op a_i,\op a_k^\dagger\op a_l]
&=\op a_i\op a_k^\dagger\op a_l-\op a_k^\dagger\op a_l\op a_i\\
&=\op a_i\op a_k^\dagger\op a_l-\mu\op a_k^\dagger\op a_i\op a_l\\
&=\op a_l\delta_{ik},\\[4pt]
[\op a_i,\op a_k^\dagger\op a_l^\dagger\op a_m\op a_n]
&=\op a_l^\dagger\op a_m\op a_n\delta_{ik}
+\op a_k^\dagger\op a_n\op a_m\delta_{il},
\end{align*}
we find
\begin{align*}
\ii\hbar\frac{\dd\op a_i}{\dd t}
&=[\op a_i,\op H]\\
&=\sum_l\op a_l\mel{a_i}{\op H_0}{a_l}
+\sum_{l,m,n}\op a_l^\dagger\op a_m\op a_n
\mel{a_i,a_l}{\op V}{a_m,a_n}.
\end{align*}
If $\op V$ has a definite value in the $A$ representation, namely
\[\op V\ket{a_m,a_n}=V_{mn}\ket{a_m,a_n},\]
then
\[\ii\hbar\frac{\dd\op a_i}{\dd t}
=\sum_l\mel{a_i}{\op H_0}{a_l}\op a_l
+\sum_lV_{il}\op a_l^\dagger\op a_l\op a_i.\]
If the representation $A$ is the continuous position-and-spin representation,
\begin{align*}
\ii\hbar\frac{\partial}{\partial t}\op\psi_\sigma(\vect r,t)
&=\sum_{\sigma'}\int\dd^3r'\,
\mel{\vect r,\sigma}{\op H_0}{\vect r',\sigma'}
\op\psi_{\sigma'}(\vect r',t)\\
&\quad+\sum_{\sigma'}\int\dd^3r'\,
V(\vect r,\vect r')
\op\psi_{\sigma'}^\dagger(\vect r',t)
\op\psi_{\sigma'}(\vect r',t)
\op\psi_\sigma(\vect r,t).
\end{align*}
If $\op H_0$ is a function only of position and momentum and is spin independent, then
\[\mel{\vect r,\sigma}{\op H_0}{\vect r',\sigma'}
=H_0(\vect r,-\ii\hbar\vect\nabla)
\delta^{(3)}(\vect r-\vect r')\delta_{\sigma\sigma'},\]
and hence
\begin{align*}
\ii\hbar\frac{\partial}{\partial t}\op\psi_\sigma(\vect r,t)
&=H_0(\vect r,-\ii\hbar\vect\nabla)\op\psi_\sigma(\vect r,t)\\
&\quad+\sum_{\sigma'}\int\dd^3r'\,
V(\vect r,\vect r')
\op\psi_{\sigma'}^\dagger(\vect r',t)
\op\psi_{\sigma'}(\vect r',t)
\op\psi_\sigma(\vect r,t).
\end{align*}
When $\op V=0$, the Heisenberg-picture annihilation operator $\op\psi_\sigma(\vect r,t)$ for the many-particle system satisfies
\[\ii\hbar\frac{\partial}{\partial t}\op\psi_\sigma(\vect r,t)
=H_0(\vect r,-\ii\hbar\vect\nabla)\op\psi_\sigma(\vect r,t).\]
This has the same form as the Schr\"odinger-picture equation for the single-particle wavefunction $\varphi_\sigma(\vect r,t)$:
\[\ii\hbar\frac{\partial}{\partial t}\varphi_\sigma(\vect r,t)
=H_0(\vect r,-\ii\hbar\vect\nabla)\varphi_\sigma(\vect r,t).\]
This shows that, if $\psi_\sigma(\vect r,t)$ is regarded as a Schr\"odinger wavefunction (a classical field), canonical quantization of this field, by imposing the commutation relations of the field operators, gives a quantum field $\op\psi_\sigma(\vect r,t)$ whose equation of motion is the same as that of the classical field $\psi_\sigma(\vect r,t)$. Therefore the procedure above, in which Fock space is introduced to treat a many-particle problem, can be regarded as \emph{second quantization}: in quantum mechanics an observable is quantized and becomes an operator, while the wavefunction remains a classical function; on entering the particle-number representation it is equivalent to quantizing the wavefunction as well, turning it into a quantum field.